\begin{proposition}
\label{proposition-finite-flat-equivalence}
Let $S$ be a scheme.
Let $(U, R, s, t, c)$ be a groupoid scheme over $S$.
Assume
\begin{enumerate}
\item $U = \Spec(A)$, and $R = \Spec(B)$ are affine,
\item $s, t : R \to U$ finite locally free, and
\item $j = (t, s)$ is an equivalence relation.
\end{enumerate}
In this case, let $C \subset A$ be as in
(\ref{equation-invariants}). Then $U \to M = \Spec(C)$
is finite locally free and $R = U \times_M U$.
Moreover, $M$ represents the quotient sheaf $U/R$
in the fppf topology (see Definition \ref{definition-quotient-sheaf}).
\end{proposition}

\begin{proof}
During this proof we use the notation $s, t : A \to B$
instead of the notation $s^\sharp, t^\sharp$.
By Lemma \ref{lemma-criterion-quotient-representable}
it suffices to show that $C \to A$ is finite locally free
and that the map
$$
t \otimes s : A \otimes_C A \longrightarrow B
$$
is an isomorphism. First, note that $j$ is a monomorphism, and
also finite (since already $s$ and $t$ are finite). Hence we see
that $j$ is a closed immersion by
Morphisms, Lemma \ref{morphisms-lemma-finite-monomorphism-closed}.
Hence $A \otimes_C A \to B$ is surjective.

\medskip\noindent
We will perform base change by flat ring maps $C \to C'$ as in
Lemma \ref{lemma-invariants-base-change}, and we will use that
formation of invariants commutes with flat base change, see
part (3) of the lemma cited.
We will show below that for every prime $\mathfrak p \subset C$, there exists
a local flat ring map $C_{\mathfrak p} \to C_{\mathfrak p}'$
such that the result holds after a base change to $C_{\mathfrak p}'$.
This implies immediately
that $A \otimes_C A \to B$ is injective (use
Algebra, Lemma \ref{algebra-lemma-characterize-zero-local}).
It also implies that $C \to A$ is flat, by combining
Algebra, Lemmas \ref{algebra-lemma-local-flat-ff},
\ref{algebra-lemma-flat-localization}, and
\ref{algebra-lemma-flatness-descends}. Then since $U \to \Spec(C)$
is surjective also (Lemma \ref{lemma-points}) we conclude that $C \to A$
is faithfully flat. Then the isomorphism $B \cong A \otimes_C A$
implies that $A$ is a finitely presented $C$-module, see
Algebra, Lemma \ref{algebra-lemma-descend-properties-modules}.
Hence $A$ is finite locally free over $C$, see
Algebra, Lemma \ref{algebra-lemma-finite-projective}.

\medskip\noindent
By Lemma \ref{lemma-finite-locally-free-disjoint-free}
we know that $A$ is a finite
product of rings $A_r$ and $B$ is a finite product of rings $B_r$
such that the groupoid scheme decomposes accordingly (see the proof
of Lemma \ref{lemma-integral-over-invariants}).
Then also $C$ is a product of rings $C_r$ and
correspondingly $C'$ decomposes as a product. Hence we may and do
assume that the ring maps $s, t : A \to B$ are finite
locally free of a fixed rank $r$.

\medskip\noindent
The local ring maps $C_{\mathfrak p} \to C_{\mathfrak p}'$ we are going
to use are any local flat ring maps such that the residue field of
$C_{\mathfrak p}'$ is infinite.
By Algebra, Lemma \ref{algebra-lemma-flat-local-given-residue-field}
such local ring maps exist.

\medskip\noindent
Assume $C$ is a local ring with maximal ideal $\mathfrak m$ and
infinite residue field, and assume that $s, t : A \to B$ is
finite locally free of constant rank $r > 0$.
Since $C \subset A$ is integral (Lemma \ref{lemma-integral-over-invariants})
all primes lying over $\mathfrak m$ are maximal, and all maximal
ideals of $A$ lie over $\mathfrak m$. Similarly for $C \subset B$.
Pick a maximal ideal $\mathfrak m'$
of $A$ lying over $\mathfrak m$ (exists by Lemma \ref{lemma-points}).
Since $t : A \to B$ is finite locally free there exist at most finitely
many maximal ideals of $B$ lying over $\mathfrak m'$. Hence we conclude
(by Lemma \ref{lemma-points} again)
that $A$ has finitely many maximal ideals, i.e.,
$A$ is semi-local. This in turn implies that $B$ is semi-local as
well. OK, and now, because $t \otimes s : A \otimes_C A \to B$ is surjective,
we can apply
Algebra, Lemma \ref{algebra-lemma-semi-local-module-basis-in-submodule}
to the ring map $C \to A$, the $A$-module $M = B$ (seen as an $A$-module
via $t$) and the $C$-submodule $s(A) \subset B$. This lemma implies that there
exist $x_1, \ldots, x_r \in A$ such that $M$ is free over $A$
on the basis $s(x_1), \ldots, s(x_r)$. Hence we conclude that $C \to A$
is finite free and $B \cong A \otimes_C A$ by applying
Lemma \ref{lemma-basis}.
\end{proof}